%% ma: L11-B16-0069
%% lop: 11
%% chuong: 5
%% bai: 16
%% dang: 11.16.3
%% loai: TLN
%% muc: VD
%% dap_an: "12"
%% nguon: "tai-lieu-hoc-tap-mon-toan-lop-11-theo-chuong-trinh-giao-duc-pho-thong-2018.pdf (Lớp Toán Thầy Duy, Phạm Lê Duy, THPT Chu Văn An), Chương 3 Bài 2, Phần 3, Câu 7 (tr.381)"
%% kien_thuc: "Lớp 11 Bài 16 (giới hạn của hàm số)"
%% trang_thai: chinh-thuc
%% ghi_chu: "Đã duyệt 03/11/2026."
\begin{ex}
Cho \[\lim\limits_{x\to-\infty}\dfrac{(3+2x^2)^5(3x^3-4)^2}{(5-x^3)^2(1-2x^2)^3}=\lim\limits_{x\to-\infty}x^4\cdot\dfrac{\left(\dfrac3{x^2}+a\right)^5\left(b-\dfrac4{x^3}\right)^2}{\left(\dfrac5{x^3}-c\right)^2\left(\dfrac1{x^2}-d\right)^3}\] với $a,b,c,d$ là các số tự nhiên. Tính tích các số $a,b,c,d$.
\shortans{12}
\loigiai{\sloppy $(3+2x^2)^5=x^{10}\left(\dfrac3{x^2}+2\right)^5\Rightarrow a=2$; $(3x^3-4)^2=x^6\left(3-\dfrac4{x^3}\right)^2\Rightarrow b=3$; $(5-x^3)^2=x^6\left(\dfrac5{x^3}-1\right)^2\Rightarrow c=1$; $(1-2x^2)^3=x^6\left(\dfrac1{x^2}-2\right)^3\Rightarrow d=2$ (và $x^{16}:x^{12}=x^4$). Tích $=2\cdot3\cdot1\cdot2=12$.}
\end{ex}
